% !TEX root = ../main.tex
% !TEX program = lualatex

\documentclass[../main.tex]{subfiles}
\IfSubfilesClassLoaded{\externaldocument{\subfix{../build/main}}}{}

\begin{document}
\IfSubfilesClassLoaded{\chapter{Production, Quality, and Manufacturing}}{}

\section{Production, Quality, and Manufacturing}

\subsection{Summary}
\begin{description}
  \item[\textbf{Production variables.}] The 5Ms and environment give a board-speed way to isolate production variation and root causes~\autocite{edo-3-7-2-production-quality-manufacturing-2026}.
  \item[\textbf{Process capability.}] A tolerance states what is acceptable; process capability and statistical control show whether the production process can repeatedly deliver it at rate~\autocite{edo-3-7-2-production-quality-manufacturing-2026,DoD-Producibility-Manufacturability-Guide-2024}.
  \item[\textbf{Quality answer.}] Quality is built into the process through design, supplier controls, manufacturing management, process proofing, and corrective action, not inspected into the product at the end~\autocite{edo-3-7-2-production-quality-manufacturing-2026,DAU-Manufacturing-Quality}.
\end{description}

\subsection{Production system variables}

Coursebook 3.7.2 ties production performance to the variables that create
process variation.  The standard board framing is the ``5Ms'': material,
machinery, method, manpower, and measurement.  Some quality tools add
environment, or milieu, as a sixth source of variation
\autocite{edo-3-7-2-production-quality-manufacturing-2026}.

\begin{longtblr}[
  caption = {Production variables and the EDO question to ask.},
  label = {tab:production-variables},
]{
  width = \linewidth,
  colspec = {X[0.7,l] X[1.6,l]},
  rowhead = 1,
  hlines,
  vlines,
}
\textbf{Variable} & \textbf{EDO question} \\
Material & Are incoming materials, parts, and supplier processes consistent enough for the required product? \\
Machinery & Are equipment, tooling, fixtures, and calibration capable at the planned rate? \\
Method & Are work instructions, routing, sequencing, and process controls mature and repeatable? \\
Manpower & Are operators trained, certified, supervised, and available? \\
Measurement & Can inspections and tests detect the variation that matters? \\
Environment & Do temperature, humidity, cleanliness, safety, or other workplace conditions affect output? \\
\end{longtblr}

\subsection{Product and process structures}

Production structures move along a variety-rate continuum.  Job shops support
high variety and low volume; batch production groups similar work; assembly
lines increase flow for more stable products; continuous-flow processes support
very high volume with highly standardized products.  The acquisition question
is whether the product design, production rate, and industrial base match the
selected process structure
\autocite{edo-3-7-2-production-quality-manufacturing-2026}.

\subsection{Variation, capability, and statistical control}

Process proofing and variation reduction precede confidence in rate production.
The program manager should distinguish tolerance from capability: a drawing
tolerance says what is acceptable, while process capability shows whether the
production process can repeatedly produce within that tolerance.  Statistical
process control is useful only when measurements are tied to the critical
process characteristics and acted on before defects escape
\autocite{edo-3-7-2-production-quality-manufacturing-2026}.

\subsection{Lean, ISO 9001, and cost of quality}

\begin{description}
  \item[\textbf{Lean.}] Lean manufacturing is a people-oriented business philosophy that focuses on value, value stream, flow, pull, and perfection.
  \item[\textbf{ISO 9001.}] DoD can accept an ISO 9001 quality system as evidence of disciplined quality management for complex or critical systems, but the program manager should not treat registration as the only possible quality approach.
  \item[\textbf{Cost of quality.}] Prevention and appraisal cost are deliberate investments to avoid internal failure, external failure, rework, warranty, and fleet-impact costs.
\end{description}
These coursebook quality concepts are linked because they all push the program
manager toward prevention, stable process control, and supplier-quality evidence
instead of late defect discovery
\autocite{edo-3-7-2-production-quality-manufacturing-2026}.

\subsection{Problem-solving tools}

Use ``5 Whys'' only after the problem statement is accurate and the answers are
complete, unbiased, and honest.  Use multivoting to reduce a long list of
candidate causes or improvements to the few priorities the team will actually
work.  Use cause-and-effect, or fishbone, diagrams to organize possible causes
by production variable before choosing corrective action
\autocite{edo-3-7-2-production-quality-manufacturing-2026}.

\subsection{Escaped Defects and Corrective Action}

When defects escape, do not answer with inspection alone.  Contain the lot,
protect fleet/user safety, identify root cause, correct the process, verify
effectiveness, assess supplier and configuration impacts, and feed lessons into
cost, schedule, and risk decisions
\autocite{edo-3-7-2-production-quality-manufacturing-2026,DoD-Producibility-Manufacturability-Guide-2024,DAU-Manufacturing-Quality}.

\paragraph{Board cue.}
Quality is built into the process, not inspected into the product.  Tie every
production-quality answer to variation, process capability, control, and cost
of quality.

\subsection{Quick Review}
\begin{itemize}
  \item What do the 5Ms help you diagnose? \textbf{Answer:} Material, machinery, method, manpower, and measurement sources of process variation; environment is often the sixth category.
  \item What is the difference between tolerance and capability? \textbf{Answer:} Tolerance is the drawing limit; capability is whether the process can repeatedly stay inside that limit.
  \item What is the defect-response trap? \textbf{Answer:} Adding inspection without containment, root-cause correction, effectiveness verification, and supplier/configuration/risk follow-through.
\end{itemize}

\IfSubfilesClassLoaded{
  \RenewDocumentCommand{\entryname}{}{\textbf{\color{Modern} Acronym}}
  \RenewDocumentCommand{\descriptionname}{}{\textbf{\color{Modern} Definition}}
  \printnoidxglossary[
    type=\acronymtype,
    title=Acronyms,
    style=long-booktabs
  ]}{}
\end{document}
